\documentclass[11pt,a4paper]{article}

\usepackage[a4paper,margin=1in]{geometry}
\usepackage{amsmath,amssymb,amsfonts,amsthm}
\usepackage{hyperref}
\usepackage{xcolor}
\usepackage[numbers,sort&compress]{natbib}

\hypersetup{colorlinks=true,linkcolor=blue,citecolor=blue,urlcolor=blue}

\newtheorem{theorem}{Theorem}
\newtheorem{proposition}[theorem]{Proposition}
\newtheorem{lemma}[theorem]{Lemma}
\newtheorem{corollary}[theorem]{Corollary}

\title{An exact closed-form for the Heinzle--Uggla average of\\
       the dominant Kasner exponent on the Mixmaster attractor}

\author{[Author names to be filled in prior to submission]}

\date{2026-05-03}

\begin{document}
\maketitle

\begin{abstract}
For the vacuum Bianchi~IX (Mixmaster) cosmology under the Belinski--Khalatnikov--%
Lifshitz approach to the past singularity, the discrete era-return map
$u\mapsto 1/\{1/u\}$ on $u\in(1,\infty)$ is conjugate to the Gauss
continued-fraction map by $x=1/u$, and the unique absolutely-continuous invariant
measure is the pushforward Gauss density $d\mu_{\rm HU}(u)=du/(u(u+1)\ln 2)$.
We give a five-line elementary proof that the time-average of the dominant Kasner
exponent at era-start under~$\mu_{\rm HU}$ admits the exact closed form
\[
   \langle p_3\rangle_{\mu_{\rm HU}}
   \;=\;\frac{\sqrt{3}\,\pi}{9\,\ln 2}
   \;=\;0.872253116\ldots
\]
The integral $\int_1^\infty p_3(u)\,du/(u(u+1))$ reduces to $\int_1^\infty
du/(1+u+u^2)=\sqrt{3}\pi/9$ by completing the square. We also provide
companion closed forms for $\langle p_1\rangle$ and $\langle p_2\rangle$
satisfying $p_1+p_2+p_3=1$, and verify them numerically at $200$-digit
precision with~\textsc{mpmath}.
\end{abstract}

% ---------------------------------------------------------------------------
\section{Setup and the BKL era-return map}
\label{sec:setup}
% ---------------------------------------------------------------------------

The canonical Kasner exponents in the standard $u$-parametrization are
\begin{equation}
  p_1(u)=\frac{-u}{1+u+u^2},\qquad
  p_2(u)=\frac{1+u}{1+u+u^2},\qquad
  p_3(u)=\frac{u(1+u)}{1+u+u^2},
\label{eq:kasner}
\end{equation}
with $u\in[1,\infty)$.  They satisfy
$p_1+p_2+p_3=1$ and $p_1^2+p_2^2+p_3^2=1$, with ordering
$p_1\le 0\le p_2\le p_3\le 1$ (Kasner identities, see e.g.\
\cite{Wainwright1997,LandauLifshitz1980}).

In the Belinski--Khalatnikov--Lifshitz (BKL) approach to the spacelike
singularity \cite{BKL1970,LandauLifshitz1980}, the dynamics of vacuum
Bianchi~IX in the asymptotic regime $t\to 0^+$ is governed by an oscillatory
sequence of Kasner epochs (Mixmaster epochs).  An \emph{era} is a maximal
sequence of within-era replacements; starting from $u_0=k+\xi$ with
$k=\lfloor u_0\rfloor\ge 2$ and $\xi=\{u_0\}\in(0,1)$, the era visits
$u_0,\,u_0-1,\ldots,u_0-(k-1)=1+\xi\in(1,2)$, and then the bounce
$1+\xi\mapsto 1/\xi$ produces the next era's start
\begin{equation}
  u_1\;=\;1/\xi\;=\;1/\{u_0\}.
\label{eq:erareturn}
\end{equation}
The era-return map $u_0\mapsto u_1$ on $(1,\infty)$ is conjugate
to the Gauss continued-fraction map $T(x)=\{1/x\}$ on $(0,1)$
via $x=1/u$.

% ---------------------------------------------------------------------------
\section{The Heinzle--Uggla measure}
% ---------------------------------------------------------------------------

The Gauss measure $d\nu(x)=dx/((1+x)\ln 2)$ is the unique
absolutely-continuous invariant measure of~$T$ (cf.\ \cite{Gauss1812,
Khinchin1964,Iosifescu2002}).  Pushing forward via $x=1/u$
gives, for $u\in(1,\infty)$,
\begin{equation}
  d\mu_{\rm HU}(u)\;=\;\frac{du}{u(u+1)\,\ln 2}.
\label{eq:HUmeasure}
\end{equation}
We refer to $\mu_{\rm HU}$ as the \emph{Heinzle--Uggla measure} after
Heinzle--Uggla \cite{HeinzleUggla2009}, who established the Mixmaster
attractor theorem on full $\mu$-measure for vacuum Bianchi~IX initial data.
Earlier rigorous statements of the era-statistics under~$\mu_{\rm HU}$
appear in \cite{KLKSS1985}.

% ---------------------------------------------------------------------------
\section{Closed-form evaluation of $\langle p_3\rangle_{\mu_{\rm HU}}$}
\label{sec:main}
% ---------------------------------------------------------------------------

\begin{theorem}[Exact mean of the dominant Kasner exponent]
\label{thm:main}
Under the Heinzle--Uggla measure \eqref{eq:HUmeasure},
\begin{equation}
  \langle p_3\rangle_{\mu_{\rm HU}}
  \;=\;\int_1^\infty p_3(u)\,d\mu_{\rm HU}(u)
  \;=\;\frac{\sqrt{3}\,\pi}{9\,\ln 2}
  \;=\;0.872253116\ldots
\label{eq:main}
\end{equation}
\end{theorem}

\begin{proof}
The integrand simplifies algebraically:
\[
  \frac{p_3(u)}{u(u+1)}
  \;=\;\frac{u(1+u)}{(1+u+u^2)\,u(u+1)}
  \;=\;\frac{1}{1+u+u^2}.
\]
Completing the square, $1+u+u^2=(u+1/2)^2+3/4$, so
\[
  \int\frac{du}{1+u+u^2}
  \;=\;\frac{2}{\sqrt{3}}\,\arctan\!\left(\frac{2u+1}{\sqrt{3}}\right)+C.
\]
Evaluating from $1$ to $\infty$,
\[
  \int_1^\infty\!\!\frac{du}{1+u+u^2}
  \;=\;\frac{2}{\sqrt{3}}\Bigl(\frac{\pi}{2}-\arctan\sqrt{3}\Bigr)
  \;=\;\frac{2}{\sqrt{3}}\Bigl(\frac{\pi}{2}-\frac{\pi}{3}\Bigr)
  \;=\;\frac{2}{\sqrt{3}}\cdot\frac{\pi}{6}
  \;=\;\frac{\pi}{3\sqrt{3}}
  \;=\;\frac{\sqrt{3}\,\pi}{9}.
\]
Dividing by $\ln 2$ gives \eqref{eq:main}.
\end{proof}

% ---------------------------------------------------------------------------
\section{Companion closed forms for $\langle p_1\rangle$ and $\langle p_2\rangle$}
% ---------------------------------------------------------------------------

The same partial-fraction technique yields
\begin{equation}
  \langle p_1\rangle_{\mu_{\rm HU}}\;=\;
  \frac{\ln 2-\tfrac{1}{2}\ln 3-\tfrac{\sqrt{3}\pi}{18}}{\ln 2}
  \;=\;-0.2286058\ldots,\qquad
  \langle p_2\rangle_{\mu_{\rm HU}}\;=\;
  \frac{\tfrac{1}{2}\ln 3-\tfrac{\sqrt{3}\pi}{18}}{\ln 2}
  \;=\;0.3563523\ldots
\label{eq:p1p2}
\end{equation}
By the Kasner identity $p_1+p_2+p_3=1$,
$\langle p_1\rangle+\langle p_2\rangle+\langle p_3\rangle=1$ holds
identically; we have verified the cancellation symbolically with
\textsc{sympy}.

\begin{proposition}[Numerical verification at $200$-digit precision]
At \textsc{mpmath} precision $\mathrm{dps}=200$, all three closed forms
\eqref{eq:main}--\eqref{eq:p1p2} agree with the directly evaluated integrals
to absolute error $0$ in the last digit, and
$\langle p_1\rangle+\langle p_2\rangle+\langle p_3\rangle-1=0$ to
$200$ digits.
\end{proposition}

% ---------------------------------------------------------------------------
\section{Remarks on related ergodic averages}
% ---------------------------------------------------------------------------

Two distinct natural ergodic measures on the Mixmaster phase space coexist
and should not be conflated.

\medskip
\noindent\textbf{(i) Era-return measure (this work)}: $\mu_{\rm HU}$ on
$(1,\infty)$, driven by the discrete era-return map \eqref{eq:erareturn}.
$\langle p_3\rangle_{\mu_{\rm HU}}=\sqrt{3}\pi/(9\ln 2)\approx 0.87225$.

\smallskip
\noindent\textbf{(ii) BKL-step-weighted measure}: each Kasner replacement
within an era is weighted by its proper-time duration~$|\ln u_n|$. By
Khinchin's theorem, the expected era length $\langle L\rangle=\int_0^1
\lfloor 1/x\rfloor\,d\nu(x)$ is divergent
\cite{Khinchin1964}, so step-weighted Birkhoff averages are
cutoff-dependent and \emph{not} a proper ergodic invariant in the usual sense;
careful regularization is required.

\medskip
The conflation of (i) and (ii) has been a source of numerical-vs-analytic
disagreement in some recent computations of Mixmaster averages.
A naive mismatch of the era-return map (replacing the conjugacy
$u_{n+1}=1/\{u_n\}$ by the off-by-one shift $u_{n+1}=1+1/\{u_n\}$)
samples instead the measure $du/((u-1)u\ln 2)$ on $(2,\infty)$, leading to
the spurious value $\langle p_3\rangle\approx 0.93578$, which we have
also evaluated symbolically and reproduced numerically.

% ---------------------------------------------------------------------------
\section*{Acknowledgements}
% ---------------------------------------------------------------------------

The closed-form evaluation in Theorem~\ref{thm:main} was independently
obtained and verified by three large-language-model agents (Sonnet, Mistral
large-latest, Opus 4.7) operating with \textsc{sympy} and
\textsc{mpmath} at $200$-digit precision; the elementary proof presented in
Section~\ref{sec:main} was confirmed against published era-statistics
results in \cite{KLKSS1985}.

% ---------------------------------------------------------------------------
\begin{thebibliography}{99}

\bibitem{BKL1970}
V.~A.~Belinskii, I.~M.~Khalatnikov, and E.~M.~Lifshitz,
\emph{Adv.\ Phys.}\ \textbf{19} (1970) 525--573.

\bibitem{LandauLifshitz1980}
L.~D.~Landau and E.~M.~Lifshitz,
\emph{The Classical Theory of Fields}, 4th rev.\ Engl.\ ed.,
Pergamon Press (1980), \S119.

\bibitem{Gauss1812}
C.~F.~Gauss, letter to Laplace, 30 January 1812; see
\cite{Iosifescu2002}~Ch.~1.

\bibitem{Khinchin1964}
A.~Ya.~Khinchin, \emph{Continued Fractions} (transl.\ from Russian),
University of Chicago Press (1964).

\bibitem{Iosifescu2002}
M.~Iosifescu and C.~Kraaikamp, \emph{Metrical Theory of Continued Fractions},
Mathematics and Its Applications \textbf{547}, Springer (2002).

\bibitem{KLKSS1985}
I.~M.~Khalatnikov, E.~M.~Lifshitz, K.~M.~Khanin, L.~N.~Shchur, and Ya.~G.~Sinai,
``On the stochasticity in relativistic cosmology,''
\emph{J.\ Stat.\ Phys.}\ \textbf{38} (1985) 97--114,
DOI \href{https://doi.org/10.1007/BF01017851}{10.1007/BF01017851}.

\bibitem{HeinzleUggla2009}
J.~M.~Heinzle and C.~Uggla,
``A new proof of the Bianchi type IX attractor theorem,''
\emph{Class.\ Quantum Grav.}\ \textbf{26} (2009) 075015,
\href{http://arxiv.org/abs/0901.0806}{arXiv:0901.0806}.

\bibitem{Wainwright1997}
J.~Wainwright and G.~F.~R.~Ellis (eds.),
\emph{Dynamical Systems in Cosmology}, Cambridge University Press (1997).

\end{thebibliography}

\end{document}
